% ============================================================
% ladder_coil_circuit.tex — coil circuit after retrofit.
%
% 3-wire seal-in with the retrofit interposed between the seal-in
% network and the overload / coil:
%
%   L1 -> STOP -> [START || M1 aux] -> TSTAT -> KA -> OL -> M1 -> L2
%
% Fail-safe by construction: every element between the rails is
% either NC-that-opens-on-fault (STOP, TSTAT, OL) or NO-that-must-
% be-actively-held-closed (START/seal, KA). Any wire break, sensor
% failure, or loss of ESP32 supply drops the coil.
%
% This is the primary audit artifact for an electrician.
%
% Regenerate with:  make ladder_coil_circuit.svg
% ============================================================
\def\pgfsysdriver{pgfsys-dvisvgm.def}
\documentclass[border=10pt,dvisvgm]{standalone}
% ============================================================
% tsstyle.tex — shared preamble for every sheet in this directory.
%
% Inputted by each sheet after \documentclass. Sheets stay pure
% geometry; everything about how the geometry *looks* lives here,
% so a style change is a one-file edit.
% ============================================================

\usepackage[T1]{fontenc}
\usepackage[utf8]{inputenc}
\usepackage[scaled=0.92]{helvet}
\renewcommand{\familydefault}{\sfdefault}

\usepackage[american,siunitx,RPvoltages]{circuitikz}
\usetikzlibrary{calc,positioning,fit,backgrounds,arrows.meta,decorations.pathreplacing}

% ------------------------------------------------------------
% Palette. Deliberately small: ink for conductors, mute for
% annotation, and exactly two semantic colours — mains (danger)
% and SELV (safe). Colour carries meaning on these sheets; it is
% never decoration.
% ------------------------------------------------------------
\definecolor{tsink}{HTML}{16181D}    % conductors, symbol bodies
\definecolor{tsmute}{HTML}{6E7178}   % secondary annotation
\definecolor{tsrule}{HTML}{C7CAD1}   % boxes, leader lines
\definecolor{tsmains}{HTML}{B3261E}  % mains potential — danger
\definecolor{tsselv}{HTML}{15607A}   % SELV / low-voltage rail
\definecolor{tssafe}{HTML}{1E6B45}   % fail-safe / passive layer
\definecolor{tspanel}{HTML}{F5F6F7}  % block fills
\definecolor{tspanelb}{HTML}{E4E7EB} % block fills, alternate

% ------------------------------------------------------------
% Global circuitikz sizing. Bipoles are scaled down from the
% default because these sheets carry a lot of elements per unit
% area; the default American resistor is comically large next to
% an IC block.
% ------------------------------------------------------------
\ctikzset{
  bipoles/length=0.85cm,
  resistors/scale=0.85,
  capacitors/scale=0.8,
  label/align=straight,
  font=\footnotesize,
}

\tikzset{
  % Conductors -------------------------------------------------
  wire/.style        = {tsink, line width=0.9pt, line cap=round},
  buswire/.style     = {tsink, line width=1.5pt, line cap=round},
  mainswire/.style   = {tsmains, line width=1.5pt, line cap=round},
  selvwire/.style    = {tsselv, line width=1.2pt, line cap=round},
  % Functional blocks -----------------------------------------
  block/.style       = {draw=tsink, line width=1pt, fill=tspanel,
                        rounded corners=1.5pt, align=center,
                        inner sep=6pt, font=\small\bfseries},
  subblock/.style    = {draw=tsrule, line width=0.7pt, fill=white,
                        rounded corners=1.5pt, align=center,
                        inner sep=4pt, font=\scriptsize},
  % Text ------------------------------------------------------
  pinlbl/.style      = {font=\scriptsize, tsink, inner sep=2pt},
  siglbl/.style      = {font=\scriptsize\itshape, tsmute, inner sep=1.5pt},
  netlbl/.style      = {font=\scriptsize\bfseries, tsink, inner sep=2pt},
  note/.style        = {font=\scriptsize, tsmute, align=left},
  danger/.style      = {font=\scriptsize\bfseries, tsmains, align=left},
  % Isolation boundary ----------------------------------------
  isoline/.style     = {tsmains, line width=1pt, dash pattern=on 5pt off 3pt},
  % A wire that crosses another without connecting. The white
  % under-stroke reads as an over-pass, so a crossing can never be
  % mistaken for a junction (junctions are always a filled dot).
  overpass/.style    = {preaction={draw, white, line width=3.4pt}},
  % Sheet furniture -------------------------------------------
  titleblock/.style  = {draw=tsink, line width=1pt, fill=white,
                        align=left, inner sep=7pt, font=\scriptsize},
  legendbox/.style   = {draw=tsrule, line width=0.7pt, fill=white,
                        align=left, inner sep=7pt, font=\scriptsize},
}

% ------------------------------------------------------------
% \pinrow{<x>}{<y>}{<dir>}{<label>}  — a pin stub on a block
% edge. dir is -1 (left edge, stub runs left) or +1 (right edge).
% ------------------------------------------------------------
\newcommand{\pinstub}[5]{%
  \draw[wire] (#1,#2) -- ++(#3*0.55,0);
  \node[pinlbl, anchor=#4] at (#1+#3*0.10,#2+0.14) {#5};
}

% ------------------------------------------------------------
% \titleblocknode{<x>}{<y>}{<sheet>}{<subtitle>}{<source>}
% Standard title block, anchored at its south east corner.
% ------------------------------------------------------------
\newcommand{\sheettitle}[5]{%
  \node[titleblock, anchor=south east] at (#1,#2) {%
    {\normalsize\bfseries #3}\\[2pt]
    {\color{tsmute}#4}\\[4pt]
    {\color{tsmute}Table-saw thermal protection retrofit\hfill}\\
    {\color{tsmute}Generated from \texttt{#5} — do not edit the SVG}%
  };
}


% ------------------------------------------------------------
% NEMA / IEC ladder primitives. CircuiTikZ ships electrical
% bipoles, not ladder shapes, so these are drawn directly. All
% three share a half-gap of 0.16 so the rung wire segments below
% line up against them.
% ------------------------------------------------------------
\def\cgap{0.16}
\def\cbar{0.34}

% \noc{<cx>}{<cy>}  — normally-open contact   -| |-
\newcommand{\noc}[2]{%
  \draw[wire] (#1-\cgap,#2-\cbar) -- (#1-\cgap,#2+\cbar);
  \draw[wire] (#1+\cgap,#2-\cbar) -- (#1+\cgap,#2+\cbar);
}
% \ncc{<cx>}{<cy>}  — normally-closed contact -|/|-
\newcommand{\ncc}[2]{%
  \noc{#1}{#2}%
  \draw[wire] (#1-0.34,#2-0.46) -- (#1+0.34,#2+0.46);
}
% \coil{<cx>}{<cy>} — output coil  -( )-
\newcommand{\coilsym}[2]{%
  \draw[wire] (#1,#2) circle (0.36);
}
% \dlabel{<cx>}{<name>}{<detail>} — name above, detail below
\newcommand{\dlabel}[3]{%
  \node[font=\small\bfseries, tsink, anchor=south] at (#1,0.86) {#2};
  \node[note, anchor=north, align=center] at (#1,-0.86) {#3};
}

\begin{document}
\begin{circuitikz}[line width=0.9pt]

% ============================================================
% Provenance banner — which hardware this sheet is drawn for
% ============================================================
\node[font=\scriptsize\bfseries, tsmains, anchor=south west] at (0,3.6)
  {PATH B TARGET --- THE THERMOSTAT (TASK-6) AND THE GPIO-DRIVEN
   RELAY (TASK-3) ARE \textbf{NOT PURCHASED}.};
\node[font=\scriptsize, tsmute, anchor=south west] at (0,3.2)
  {The coil circuit as it can be built today is on
   \texttt{pathA\_ladder\_coil\_circuit.svg}.};

% ============================================================
% Retrofit highlight — drawn first, behind everything
% ============================================================
\begin{scope}[on background layer]
  \fill[tssafe!7, rounded corners=4pt] (11.4,-2.2) rectangle (16.2,1.85);
  \draw[tssafe!35, line width=0.8pt, rounded corners=4pt]
        (11.4,-2.2) rectangle (16.2,1.85);
\end{scope}
\node[font=\scriptsize\bfseries, tssafe, anchor=south] at (13.8,1.92)
  {ADDED BY THE RETROFIT};

% ============================================================
% Rails
% ============================================================
\draw[buswire] (0,2.6) -- (0,-4.8);
\draw[buswire] (22,2.6) -- (22,-4.8);
\node[font=\small\bfseries, anchor=south] at (0,2.7) {L1};
% The A202C is available with either a 120 V or a 240 V coil and the
% saw's has not been read off the coil label yet. Do not assert a
% number here — it changes what the control rails are and, on the
% Path A sheet, what the RF receiver's supply is tapped from.
\node[note, anchor=north west] at (0.15,2.62) {control circuit hot};
\node[font=\small\bfseries, anchor=south] at (22,2.7) {L2};
\node[note, anchor=north east] at (21.85,2.62) {control circuit return};
\node[note, anchor=east] at (-0.25,0) {Rung 1};

% ============================================================
% Rung — wire segments between elements
% ============================================================
\draw[wire] (0,0)     -- (2.64,0);
\draw[wire] (2.96,0)  -- (5.20,0);
\draw[wire] (5.20,0)  -- (7.24,0);
\draw[wire] (7.56,0)  -- (9.80,0);
\draw[wire] (9.80,0)  -- (12.24,0);
\draw[wire] (12.56,0) -- (15.04,0);
\draw[wire] (15.36,0) -- (17.84,0);
\draw[wire] (18.16,0) -- (20.24,0);
\draw[wire] (20.96,0) -- (22,0);

% ============================================================
% Seal-in branch — M1 aux in parallel with START
% ============================================================
\draw[wire] (5.20,0) -- (5.20,-3.2);
\draw[wire] (9.80,0) -- (9.80,-3.2);
\draw[wire] (5.20,-3.2) -- (7.24,-3.2);
\draw[wire] (7.56,-3.2) -- (9.80,-3.2);
\draw (5.20,0) node[circ]{};
\draw (9.80,0) node[circ]{};

% ============================================================
% Elements, left to right
% ============================================================
\ncc{2.8}{0}
\dlabel{2.8}{STOP}{NC pushbutton}

\noc{7.4}{0}
\dlabel{7.4}{START}{NO pushbutton,\\momentary}

\noc{7.4}{-3.2}
\node[font=\small\bfseries, tsink, anchor=south] at (7.4,-2.34) {M1 aux};
\node[note, anchor=north, align=center] at (7.4,-4.06)
  {NO auxiliary contact on M1.\\Carries the seal after START releases.};

% TSTAT must open ABOVE the firmware's TRIP_THRESHOLD (110 °C) so
% the ESP32 always acts first and the thermostat stays a backstop.
% ARCHITECTURE.md § TASK-6 specifies 120–130 °C; this label used to
% read 110 °C, which would have put both layers on the same number.
\ncc{12.4}{0}
\dlabel{12.4}{TSTAT}{NC bimetallic,\\opens 120--130\,\textdegree C,\\auto-reset}

\noc{15.2}{0}
\dlabel{15.2}{KA}{ESP32 relay, NO,\\opto-isolated,\\active-HIGH}

\ncc{18.0}{0}
\dlabel{18.0}{OL}{NC overload aux}

\coilsym{20.6}{0}
\dlabel{20.6}{M1}{contactor coil,\\Gould/ITE A202C,\\NEMA Size 1.\\Read the coil voltage\\off the coil label.}

% ============================================================
% Series callout under TSTAT + KA
% ============================================================
\draw[tssafe, line width=0.8pt]
  (11.9,-2.6) -- (11.9,-2.8) -- (15.7,-2.8) -- (15.7,-2.6);
\node[font=\scriptsize\bfseries, tssafe, anchor=north] at (13.8,-2.87)
  {IN SERIES — either one open drops the coil};

% ============================================================
% Legend
% ============================================================
\node[legendbox, anchor=north west, text width=10.6cm] at (0,-6.0) {%
  \textbf{Why this topology is fail-safe}\\[4pt]
  Every element between the rails is either \emph{NC that opens on
  fault} or \emph{NO that must be actively held closed}. There is no
  state in which a failure leaves the coil energised.\\[5pt]
  \begin{tabular}{@{}p{1.35cm}@{\hspace{5pt}}p{1.15cm}@{\hspace{5pt}}p{5.4cm}@{}}
    \textbf{Element} & \textbf{Type} & \textbf{Opens on}\\[2pt]
    STOP   & NC & operator command; broken wire also opens\\[1pt]
    START  & NO & --- (must be held until M1 aux seals in)\\[1pt]
    M1 aux & NO & loss of coil power breaks the seal\\[1pt]
    TSTAT  & NC & winding 120--130\,\textdegree C, or a broken lead\\[1pt]
    KA     & NO & anything that stops GPIO26 being driven HIGH\\[1pt]
    OL     & NC & sustained overcurrent\\
  \end{tabular}
};

\node[legendbox, anchor=north west, text width=10.2cm] at (11.6,-6.0) {%
  \textbf{What opens KA}\\[4pt]
  KA closes only while the ESP32 actively drives GPIO26 HIGH.
  Each of these opens it:\\[3pt]
  \begin{tabular}{@{}p{4.55cm}@{\hspace{5pt}}p{4.5cm}@{}}
    ESP32 crash / brown-out / power loss & GPIO26 floats, the
      10\,k$\Omega$ pulldown wins\\[2pt]
    Winding overtemp on the K-type & firmware drives GPIO26 LOW\\[2pt]
    Thermocouple open or shorted & firmware trips on sensor loss\\
  \end{tabular}\\[6pt]
  \textcolor{tssafe}{\textbf{SR-3}}\quad TSTAT and KA are in series.
  The passive thermostat cuts the coil with no firmware involved.\\[3pt]
  \textcolor{tssafe}{\textbf{SR-4}}\quad No autonomous restart. Once the
  coil drops the seal-in is broken, so a human must press START.
};

% ============================================================
% Title block
% ============================================================
\node[titleblock, anchor=north west, text width=6.4cm] at (0,-11.0) {%
  {\normalsize\bfseries Coil circuit — ladder}\\[2pt]
  {\color{tsmute}3-wire seal-in with the retrofit interposed}\\[5pt]
  {\color{tsmute}Table-saw thermal protection retrofit}\\
  {\color{tsmute}Harness: \texttt{../harness/mains\_and\_coil.yml}}\\
  {\color{tsmute}Source: \texttt{ladder\_coil\_circuit.tex}}%
};

\node[danger, anchor=north west, text width=9.5cm] at (11.6,-11.0) {%
  Every conductor on this sheet is at mains potential.
  Lock out the disconnect before touching the starter enclosure.};

\end{circuitikz}
\end{document}
